% section: 母関数：無限個の数を一つの対象にする

\section{母関数：無限個の数を一つの対象にする}

\subsection{なぜ数列を式に詰め込むのか}

数列は無限個の項を持つ.
項を一つずつ扱う代わりに,すべての項を係数として一つの式にまとめる.
\[
  A(x)=a_0+a_1x+a_2x^2+\cdots
\]
これを母関数という.

母関数では,$x^n$ の係数が $a_n$ である.
したがって,母関数を変形して係数を読むことができれば,数列の問題を式の問題へ移せる.

\subsection{形式的冪級数としての母関数}

ここで,母関数を通常の関数と考える必要はない.
例えば,
\[
  1+x+x^2+x^3+\cdots
\]
は,$|x|<1$ では $1/(1-x)$ に収束するが,形式的冪級数としては,係数を並べた無限列である.

形式的に,
\[
  (1-x)(1+x+x^2+\cdots)=1
\]
と係数を比較できる.
収束するかどうかを考えなくても,漸化式の係数操作が可能になる.

一方,実際に数値を代入して関数として使う場合には,収束範囲を確認する必要がある.
同じ記号でも,形式的な議論と解析的な議論を区別しなければならない.

\subsection{線形漸化式を母関数へ移す}

\[
  a_{n+2}=a_{n+1}+a_n
\]
とし,
\[
  A(x)=\sum_{n=0}^{\infty}a_nx^n
\]
と置く.
漸化式に $x^n$ を掛けて足すと,
\[
  \sum_{n=0}^{\infty}a_{n+2}x^n
  =
  \sum_{n=0}^{\infty}a_{n+1}x^n
  +
  \sum_{n=0}^{\infty}a_nx^n
\]
である.

左辺と右辺を $A(x)$ で表すと,
\[
  \frac{A(x)-a_0-a_1x}{x^2}
  =
  \frac{A(x)-a_0}{x}+A(x)
\]
となる.
これを整理して,
\[
  A(x)=\frac{a_0+(a_1-a_0)x}{1-x-x^2}
\]
を得る.

分母の多項式 $1-x-x^2$ は,漸化式の特性方程式と同じ情報を持つ.
母関数は,特性方程式を別の形で記録しているのである.

\subsection{畳み込みが積になる}

母関数の特に重要な性質は,積の係数である.
\[
  A(x)=\sum_{n=0}^{\infty}a_nx^n,\qquad
  B(x)=\sum_{n=0}^{\infty}b_nx^n
\]
とすると,
\[
  A(x)B(x)
  =
  \sum_{n=0}^{\infty}
  \left(\sum_{k=0}^{n}a_kb_{n-k}\right)x^n
\]
となる.

内側の和
\[
  \sum_{k=0}^{n}a_kb_{n-k}
\]
を畳み込みという.
二つの選択を組み合わせる数え上げでは,畳み込みが自然に現れる.
母関数は,畳み込みを積へ変えることで,組合せ論の漸化式を代数方程式へ移す.

\subsection{カタラン数：非線形漸化式を解く}

カタラン数 $C_n$ を,
\[
  C_0=1,\qquad
  C_{n+1}=\sum_{k=0}^{n}C_kC_{n-k}
\]
で定める.
最初の項は,
\[
  1,1,2,5,14,42,\ldots
\]
である.

母関数
\[
  C(x)=\sum_{n=0}^{\infty}C_nx^n
\]
を考える.
漸化式の右辺は畳み込みなので,母関数では積になる.
\[
  C(x)-1=xC(x)^2
\]
すなわち,
\[
  xC(x)^2-C(x)+1=0
\]
である.

二次方程式として解くと,
\[
  C(x)=\frac{1\pm\sqrt{1-4x}}{2x}
\]
となる.
$x=0$ の近くで $C(x)\to C_0=1$ でなければならないので,正しい枝は,
\[
  C(x)=\frac{1-\sqrt{1-4x}}{2x}
\]
である.

もともとの漸化式は,各項がそれ以前の項の積の和で表される非線形なものだった.
それが,母関数を導入すると二次方程式になった.
母関数は線形漸化式専用の道具ではない.
「組合せの分解が畳み込みになる」場面では,非線形な漸化式にも力を発揮する.

\subsection{母関数から見る漸化式の意味}

母関数による変換は,
\[
  \text{無限個の数}
  \longrightarrow
  \text{一つの形式的な対象}
\]
である.
この変換によって,シフトは $x$ の掛け算になり,畳み込みは積になり,漸化式は代数方程式になる.

何を簡単にしたいかによって,どの対象へ移すかが決まる.
差分では変化を簡単にし,行列では連立した状態を簡単にし,母関数では無限個の項の関係を簡単にする.
